% from ega0-19-fr.tex
\oldpage[0\textsubscript{IV-1}]{105}

  \item[b)] $B$ est un $A$-module \emph{plat} et
  $B_0=B/\mathfrak mB$ (muni de la topologie quotient) est une $k$-algèbre
  \emph{formellement lisse}.
\end{enumerate}
\end{theorem}

\begin{proof}
La démonstration se fait en plusieurs étapes.

\medskip
\noindent\phantomsection\label{0.19.7.1.1-fr}\textbf{(19.7.1.1)}\enspace
Démontrons d'abord que $b)$ entraîne $a)$~; nous allons appliquer le critère
(19.4.7), avec $\mathfrak J=\mathfrak m$~; en vertu de la seconde hypothèse
dans $b)$, tout revient à montrer que $B$ est un $A$-module
\emph{formellement projectif}. L'hypothèse entraîne que pour tout $h>0$,
$B/\mathfrak m^hB$ est un $(A/\mathfrak m^h)$-module plat
($\mathbf 0_{\mathrm{III}}$, 10.2.1)~; comme les $\mathfrak m^h$ forment un
système fondamental de voisinages de 0 dans $A$ et que
$(B/\mathfrak m^hB)\otimes_{A/\mathfrak m^h}k=B_0$, on peut remplacer $A$ et
$B$ par $A/\mathfrak m^h$ et $B/\mathfrak m^hB$ respectivement, et par suite
supposer $A$ \emph{artinien} (donc discret). Comme $B_0$ est une $k$-algèbre
formellement lisse, c'est un anneau \emph{régulier} (19.6.5)~; soit
$(x_i^0)_{1\leq i\leq n}$ un système régulier de paramètres pour $B_0$
(17.1.6), et pour tout $i$, soit $x_i\in B$ tel que $x_i^0$ soit son image
dans $B_0=B/\mathfrak mB$~; comme les $x_i^0$ engendrent l'idéal maximal
$\mathfrak n_0=\mathfrak n/\mathfrak mB$ de $B_0$, les idéaux
\[
  \mathfrak Q_m=\sum_{i=1}^n(x_i^0)^mB_0 \qquad (m>0)
\]
forment un système fondamental de voisinages de 0 dans $B_0$, car
$\mathfrak Q_m$ est évidemment contenu dans $\mathfrak n_0^m$, et d'autre
part contient $\mathfrak n_0^{mn}$. Posons
$\mathfrak J_m=\sum_{i=1}^n x_i^mB$ pour tout $m>0$~; il est clair que
$\mathfrak n=\mathfrak J_1+\mathfrak mB$~; comme il existe un $h>0$ tel que
$\mathfrak m^h=0$, on a
$\mathfrak n^m\subset\mathfrak J_1^{m-h}\subset\mathfrak n^{m-h}$ pour
$m>h$, et comme on a vu que $\mathfrak J_m\supset\mathfrak J_1^{mn}$, on voit
que les $\mathfrak J_m$ forment un système fondamental de voisinages de 0 dans
$B$. Tout revient par suite à prouver que les $B/\mathfrak J_m$ sont des
$A$-modules \emph{libres}, et il revient au même de voir que ce sont des
$A$-modules \emph{plats} ($\mathbf 0_{\mathrm{III}}$, 10.1.3). Or,
l'hypothèse que $(x_i^0)$ est une suite $B_0$-régulière d'éléments de l'idéal
maximal de $B_0$ entraîne la même propriété pour la suite des $(x_i^0)^m$
($1\leq i\leq n$) pour tout $m>0$ (15.1.20)~; la conclusion résulte donc de
(15.1.16, $b)$ et $c)$).
\end{proof}

\begin{lemma}[19.7.1.2]
\label{0.19.7.1.2-fr}
Soient $A$ un anneau topologique, $B$, $C$ deux $A$-algèbres topologiques qui
sont des anneaux locaux noethériens. On suppose en outre que $C$ est complet
et que le corps résiduel $B/\mathfrak m$ de $B$ est un $A$-module de type fini.
Soit $E$ le produit tensoriel complété $B\widehat\otimes_A C$. Alors~:
\begin{enumerate}
  \item[(i)] $E$ est un anneau semi-local noethérien complet.
  \item[(ii)] L'idéal $\mathfrak mE$ est contenu dans le radical de $E$, et
  pour tout $h>0$, $E/\mathfrak m^hE$ est isomorphe à
  $(B/\mathfrak m^h)\otimes_A C$.
  \item[(iii)] Si $C$ est un $A$-module plat, $E$ est un $B$-module plat.
\end{enumerate}
\end{lemma}

\begin{proof}
Par définition, $E$ est le séparé complété du produit tensoriel $B\otimes_A C$
pour la topologie définie par les idéaux
$\operatorname{Im}(\mathfrak m^i\otimes_A C)+
\operatorname{Im}(B\otimes_A\mathfrak n'{}^j)$
($\mathbf 0_{\mathrm I}$, 7.7.5). Si l'on pose
\[
  \mathfrak r=\operatorname{Im}(\mathfrak m\otimes_A C)+
  \operatorname{Im}(B\otimes_A\mathfrak n'),
\]
on a
$\mathfrak r^{2i}\subset
\operatorname{Im}(\mathfrak m^i\otimes_A C)+
\operatorname{Im}(B\otimes_A\mathfrak n'{}^i)\subset\mathfrak r^i$, donc $E$
est aussi le séparé complété de $B\otimes_A C$ pour la topologie
$\mathfrak r$-préadique. Par hypothèse,
$(B\otimes_A C)/\mathfrak r=(B/\mathfrak m)\otimes_A(C/\mathfrak n')$ est un
$(C/\mathfrak n')$-module de type fini, donc un anneau artinien~; en outre,
$\mathfrak r/\mathfrak r^2$, étant un quotient de
\[
  ((\mathfrak m/\mathfrak m^2)\otimes_A C)
  \oplus(B\otimes_A(\mathfrak n'/\mathfrak n'{}^2)),
\]
est un $(B\otimes_A C)$-module de type fini~; appliquant
($\mathbf 0_{\mathrm I}$, 7.2.11), on voit donc que $E$ est un anneau
\emph{noethérien}~; en outre, $E/\mathfrak rE$, isomorphe à
$(B\otimes_A C)/\mathfrak r$, étant artinien, $E$ est \emph{semi-local}.
Notons maintenant que $E$, qui est isomorphe à
\[
  \varprojlim_{i,j}((B/\mathfrak m^i)\otimes_A(C/\mathfrak n'{}^j)),
\]
est aussi isomorphe, par le théorème de la double limite projective, à
\[
  \varprojlim_i\left(\varprojlim_j
  ((B/\mathfrak m^i)\otimes_A(C/\mathfrak n'{}^j))\right)~;
\]

